\documentclass[11pt,a4paper,openany,oneside]{book}



\usepackage[UTF8]{ctex}

\usepackage{amsmath,amssymb,mathdots}
\usepackage{mathrsfs}
\usepackage{enumerate}
%\usepackage{geometry}

%页眉
\usepackage{bookmark}
\usepackage{hyperref}
\usepackage{fancyhdr}
\renewcommand{\headwidth}{\textwidth}

\pagestyle{fancy}                       
\renewcommand{\headrulewidth}{0.5pt}
\fancyhf{}                                                
\fancyhead[LE,RO]{\thepage}                      
\fancyhead[RE]{\leftmark}                               
\fancyhead[LO]{\rightmark}

\fancypagestyle{plain}{
	\fancyhf{}                               
	\renewcommand{\headrulewidth}{0.5pt}                              
	\fancyhead[LE,RO]{\thepage}                         
	\fancyhead[RE]{\leftmark}                               
	\fancyhead[LO]{\rightmark} 
}

\newcommand{\headodd}[2]{\fancypagestyle{myplain}{% 
		\fancyhf{} %  
		\pagestyle{fancy}                              
		\renewcommand{\headrulewidth}{0.5pt}                                    
		\fancyhead[LE,RO]{\thepage}                        
		\fancyhead[RE]{\leftmark} 
		\hypersetup{   
			colorlinks=false,
			pdfborder={0 0 0},
		}                            
		\fancyhead[LO]{ \ref{#1} #2}  %#1表示label, #2表示节名
	}
	\thispagestyle{myplain}
}




%定义分数线

%\geometry{left=6.0cm, right=3.0cm, top=3.5cm, bottom=2.5cm}
\setlength{\parindent}{0pt}
%\renewcommand{\baselinestretch}{1.2}
\newcommand\chfrac[2]{\frac{\;#1\;}{\;#2\;}}
\makeatletter
\renewcommand*\env@matrix[1][\arraystretch]{%
	\edef\arraystretch{#1}%
	\hskip -\arraycolsep
	\let\@ifnextchar\new@ifnextchar
	\array{*\c@MaxMatrixCols c}}
\makeatother
\usepackage{arydshln}
\usepackage{galois}

%定理
\usepackage{lipsum}
\usepackage[many]{tcolorbox}

\definecolor{greentitle}{RGB}{61,170,61}
\definecolor{greentitleback}{RGB}{216,233,213}
%[
%auto counter,
%number within=section
%]
\newtcolorbox{myexample}[2][]{%
	breakable,
	enhanced,
	colback=white,
	colbacktitle=white,
	arc=0pt,
	leftrule=1.4pt,
	rightrule=0pt,
	toprule=0pt,
	bottomrule=0pt,
	titlerule=-16pt,
	colframe=greentitleback,
	fonttitle=\normalcolor,
	overlay={
		\node[
		outer sep=0pt,
		anchor=east,
		text width=2.2cm,
		minimum height=4ex,
		fill=greentitleback,
		font=\color{greentitle}\sffamily\scshape
		] at (title.west) {problem};
	},
	title=#2,
    #1
}
\newcommand\Solution{\noindent\textbf{\textsf{Solution}}\par\medskip}




%%%%%%%%%%%

\usepackage{tikz}
\usepackage{titletoc}
\usepackage{titlesec}
\usepackage{hyperref}
\usepackage{times}

\definecolor{doc}{RGB}{61,170,61}
\definecolor{myblueii}{RGB}{119,136,153}

%~~~~~~~~~~~~~~~~~Hyperref Settings~~~~~~~~~~~~~~~~~~~~~~

\hypersetup{%
	pdfborder = {0 0 0},
	colorlinks,
	citecolor=red,
	filecolor=green,
	linkcolor=myblueii,
	urlcolor=cyan!50!black!90
}

%~~~~~~~~~~~~~~~~~~~~~Fancy ToC~~~~~~~~~~~~~~~~~~~~~~

\contentsmargin{0cm}
\titlecontents{chapter}[0pc]
{\addvspace{30pt}%
	\begin{tikzpicture}[remember picture, overlay]%
	\draw[fill=myblueii,draw=myblueii, rounded corners] (-4,-.1) rectangle (-0.15,.5);%
	\pgftext[left,x=-2.7cm,y=0.2cm]{\color{white}\Large \chaptertitlename\ \thecontentslabel};%
	\end{tikzpicture}\color{myblueii}\large\bfseries}%
{}
{}
{\hspace*{6pt}\titlerule\hspace*{6pt}\large\bfseries \thecontentspage
	\begin{tikzpicture}[remember picture, overlay]
	\draw[fill=doc!25,draw=myblueii, rounded corners=0pt] (2pt,0) rectangle (6,0.1pt);
	\end{tikzpicture}}%
\titlecontents{section}[2.4pc]
{\addvspace{1pt}}
{\contentslabel[\thecontentslabel]{2.4pc}}
{}
{\hfill\small \thecontentspage}
[]
\titlecontents{subsection}[4.8pc]
{\addvspace{1.0pt}}
{\contentslabel[\thecontentslabel]{2.4pc}}
{}
{\hfill\small\thecontentspage}
[]

\makeatletter
\renewcommand{\tableofcontents}{%
	\chapter*{%
		\vspace*{-20\p@}%
		\begin{tikzpicture}[remember picture, overlay]%
		\pgftext[right,x=15cm,y=0.2cm]{\color{myblueii}\Huge \contentsname};%
		\draw[fill=myblueii,draw=myblueii, rounded corners=15pt] (13,-.75) rectangle (20,1);%
		\clip (13,-.75) rectangle (20,1);
		\pgftext[right,x=15cm,y=0.2cm]{\color{white}\Huge \contentsname};%
		\end{tikzpicture}}%
	\@starttoc{toc}}
\makeatother



 \usepackage{graphicx} 
 \newcommand{\HRule}{\rule{\linewidth}{0.5mm}} 
 
%~~~~~~~~~~~~~~~~~Main Document~~~~~~~~~~~~~~~~~~~~~~~

\begin{document}


\begin{titlepage}
	
	\begin{center}
		
		
		% Upper part of the page
		    
		
		%\textsc{\LARGE University of Beer}\\[1.5cm]
		
		%\textsc{\Large Final year project}\\[0.5cm]
		
		
		% Title
		\HRule \\[1cm]
		{ \Huge \bfseries 数学分析杂题}\\[0.5cm]
		
		\HRule \\[1.5cm]
		
		% Author and supervisor
		%\begin{minipage}{0.4\textwidth}
			%\begin{flushleft} \large
				%\emph{Author:}\\
		%		John \textsc{Smith}
		%	\end{flushleft}
		%\end{minipage}
		\begin{minipage}{0.9\textwidth}
			\begin{flushright} \large
				\LARGE\emph{整理}: 
				\textsc{\huge Chen}
			\end{flushright}
		\end{minipage}

	    \vfill
		
		
		
		%\includegraphics[width=0.048\textwidth]{./logo}%\\[6cm]
		% Bottom of the page
%{\hspace{5em}\Large\today\hspace{3em}}
	%\includegraphics[width=0.048\textwidth]{./logo}%\\[6cm]
	\end{center}
	
	\end{titlepage}


{   
	\hypersetup{hidelinks}
	\tableofcontents% 
}
\thispagestyle{empty}


\part{上册}


\chapter{级数}
\section{Demo}



\begin{myexample}
\textbf{1.} 
证明 $\sum\limits_{m=0}^{\infty}\dfrac{1}{(2m+1)^2}=\dfrac{\pi^2}{8}$.
\end{myexample}

\Solution


当 $0<x<1$ 时, 作代换 $\sqrt{1-x^2 \sin ^2\theta }=u\in[\sqrt{1-x^2},1]$, 即 $\theta=\arcsin\dfrac{\sqrt{1-u^2}}{x}$, 于是
\begin{align*}
\int_{0}^{\frac{\pi}{2}}\frac{\sin\theta }{\sqrt{1-x^2 \sin ^2\theta }}d\theta&=\int_{1}^{\sqrt{1-x^2}}\frac{-2}{x\sqrt{u^2-(1-x^2)}}du\\
&=\frac{1}{x}\int_{\sqrt{1-x^2}}^{1}\frac{1}{\sqrt{u^2-(1-x^2)}}du\\
&=\frac{1}{x}\left[\ln\left|u+\sqrt{u^2-(1-x^2)}\right|\right]_{\sqrt{1-x^2}}^{1}\\
&=\frac{1}{x}\left(\ln(1+x)-\ln\sqrt{1-x^2}\right)\\
&=\frac{1}{2x}\ln(1+x)-\frac{1}{2x}\ln(1-x)\\
&=\frac{1}{2x}\sum_{n=1}^{\infty}\frac{(-1)^{n-1}x^{n}}{n}-\frac{1}{2x}\sum_{n=1}^{\infty}\frac{(-1)^{n-1}(-x)^{n}}{n}\\
&=\sum_{n=1}^{\infty}\frac{(-1)^{n-1}x^{n-1}+x^{n-1}}{2n}\quad(\text{令}\,n=2m+1)\\
&=\sum_{m=0}^{\infty}\frac{x^{2m}}{2m+1}\quad(0<x<1).
\end{align*}
当 $x=0$ 时, 
\[
\int_{0}^{\frac{\pi}{2}}\frac{\sin\theta }{\sqrt{1-x^2 \sin ^2\theta }}d\theta=\sum_{m=0}^{\infty}\frac{x^{2m}}{2m+1}
\]
仍然成立.
考虑瑕点为 $x=1$ 的广义积分, 有
\begin{align}
\int_{0}^{1}dx\int_{0}^{\frac{\pi}{2}}\frac{\sin\theta }{\sqrt{1-x^2 \sin ^2\theta }}d\theta&=\int_{0}^{1}\sum_{m=0}^{\infty}\frac{x^{2m}}{2m+1} dx\\
&=\lim_{t\to 1^-}\int_{0}^{t}\sum_{m=0}^{\infty}\frac{x^{2m}}{2m+1} dx\\
&=\lim_{t\to 1^-}\sum_{m=0}^{\infty}\int_{0}^{t}\frac{x^{2m}}{2m+1} dx\\
&=\lim_{t\to 1^-}\sum_{m=0}^{\infty}\frac{t}{(2m+1)^2} \\
&=\sum_{m=0}^{\infty}\lim_{t\to 1^-}\frac{t}{(2m+1)^2} \\
&=\sum_{m=0}^{\infty}\frac{1}{(2m+1)^2}.
\end{align}

另一方面, 
\begin{align*}
\int_{0}^{\frac{\pi}{2}}d\theta\int_{0}^{1}\frac{\sin\theta }{\sqrt{1-x^2 \sin ^2\theta }}dx
=\int_{0}^{\frac{\pi}{2}}\left[\arcsin(x\sin\theta)\right]_0^1d\theta
=\int_{0}^{\frac{\pi}{2}}\theta d\theta=\frac{\pi^2}{8}	.
\end{align*}

这表明
\[\sum\limits_{n=1}^{\infty}\dfrac{1}{(2n-1)^2}=\dfrac{\pi^2}{8}.\]

~\\

\textbf{注 1}：由 1.2 式至 1.3 式交换了积分号与求和号, 这里逐项积分的合法性由幂级数 \[\sum\limits_{m=0}^{\infty}\dfrac{x^{2m}}{2m+1}\]
在 $[0,t]$ 上一致收敛所保证. 事实上, 该幂级数的收敛域为 $(-1,1)$. 注意到在 $t\to 1^{-}$ 时有 $t<1$, 由 Abel 第二定理可知幂级数于其收敛域中内闭一致收敛.\\
~\\
\textbf{注 2}：由 1.4 式至 1.5 式交换了极限符号与求和号. 为此, 我们可以验证以 $t$ 为变元的级数
\[
\sum_{m=0}^{\infty}\frac{t}{(2m+1)^2}
\]
在 $[0,2]$ 上一致收敛. 因为
\[
\sum_{m=0}^{\infty}\frac{t}{(2m+1)^2}\le\sum_{m=0}^{\infty}\frac{3}{(2m+1)^2}\le3+\sum_{m=1}^{\infty}\frac{3}{(2m-1)(2m+1)}<6,
\]
运用 Weierstrass 判别法即可得到证明.
~\\
~\\

\textbf{另证:}\\
我们另外给出一个证明, 作为对极限交换技巧的一种展示.\\
定义双阶乘记号
\[
(2n-1)!!=(2n-1)(2n-3)\cdots3\cdot1,
\]
\[
(2n)!!=(2n)(2n-2)\cdots4\cdot2.
\]
写出 $(1+z)^{-\frac{1}{2}}$ 在 $z=0$ 处的 Taylor 级数展开式
\begin{align*}
(1+z)^{-\frac{1}{2}}&=1+\frac{-\frac{1}{2}}{1!}z+\frac{(-\frac{1}{2})(-\frac{1}{2}-1)}{2!}z^2+\frac{(-\frac{1}{2})(-\frac{1}{2}-1)(-\frac{1}{2}-2)}{3!}z^3+\cdots\\
&=1+\sum_{n=1}^{\infty}\frac{(-\frac{1}{2})(-\frac{3}{2})\cdots(-\frac{1}{2}-n+1)}{n!}z^n\\
&=1+\sum_{n=1}^{\infty}\frac{(-1)^n(2n-1)!!}{2^n n!}z^n\\
&=1+\sum_{n=1}^{\infty}\frac{(-1)^n(2n-1)!!}{(2n)!!}z^n\quad (-1<z\le 1).
\end{align*}


令 $z=-x^2 \sin ^2\theta$, 则得到 $f(x)=\frac{\sin\theta }{\sqrt{1-x^2 \sin ^2\theta }}$ 在 $x=0$ 处的 Taylor 级数展开式
\begin{align*}
\frac{\sin\theta }{\sqrt{1-x^2 \sin ^2\theta }}&=\sin\theta\left(1+\sum_{n=1}^{\infty}\frac{(-1)^n(2n-1)!!}{(2n)!!}(-x^2 \sin ^2\theta)^n\right)\\
&=\sin\theta+\sum_{n=1}^{\infty}\frac{(2n-1)!!}{(2n)!!}x^{2n} \sin ^{2n+1}\theta\quad(-1<x<1).
\end{align*}

进一步有

\begin{align}
&\int_{0}^{\frac{\pi}{2}}d\theta\int_{0}^{1}\frac{\sin\theta }{\sqrt{1-x^2 \sin ^2\theta }}dx\\
=&\int_{0}^{\frac{\pi}{2}}d\theta\lim_{t\to 1^-}\int_{0}^{t}\left(\sin\theta+\sum_{n=1}^{\infty}\frac{(2n-1)!!}{(2n)!!}x^{2n} \sin ^{2n+1}\theta\right)dx\\
=&\int_{0}^{\frac{\pi}{2}}d\theta\lim_{t\to 1^-}\left(\int_{0}^{t}\sin\theta dx+\sum_{n=1}^{\infty}\frac{(2n-1)!!}{(2n)!!}\sin ^{2n+1}\theta\int_{0}^{t}x^{2n}dx \right)\\
=&\int_{0}^{\frac{\pi}{2}}\left(\sin\theta+\lim_{t\to 1^-}\sum_{n=1}^{\infty}\frac{(2n-1)!!}{(2n+1)(2n)!!}t^{2n+1} \sin ^{2n+1}\theta\right)d\theta\\
=&\int_{0}^{\frac{\pi}{2}}\left(\sin\theta+\sum_{n=1}^{\infty}\lim_{t\to 1^-}\frac{(2n-1)!!}{(2n+1)(2n)!!}t^{2n+1} \sin ^{2n+1}\theta\right)d\theta\\
=&\int_{0}^{\frac{\pi}{2}}\left(\sin\theta+\sum_{n=1}^{\infty}\frac{(2n-1)!!}{(2n+1)(2n)!!}\sin ^{2n+1}\theta\right)d\theta\\
=&\int_{0}^{\frac{\pi}{2}}\sin\theta d\theta+\sum_{n=1}^{\infty}\frac{(2n-1)!!}{(2n+1)(2n)!!} \int_{0}^{\frac{\pi}{2}}\sin ^{2n+1}\theta d\theta\\
=&1+\sum_{n=1}^{\infty}\frac{(2n-1)!!}{(2n+1)(2n)!!} \frac{(2n)!!}{(2n+1)!!}\\
=&1+\sum_{n=1}^{\infty}\frac{1}{(2n+1)^2}\\
=&\sum\limits_{m=0}^{\infty}\dfrac{1}{(2m+1)^2}.
\end{align}

接着证明上述所有交换极限的行为都是合理的.

\begin{enumerate}[I.]
\item (1.8) 至 (1.9). 运用 Cauchy 根值判别法以及 Wallis 公式, 可得 $|x|<1$ 时,
\[
\rho=1\Bigg/\,\overline{\lim\limits_{n\to\infty}}\sqrt[^{^{^n}}]{\frac{(2n-1)!!}{(2n)!!} \sin ^{2n+1}\theta}=1\Bigg/\,\overline{\lim\limits_{n\to\infty}}\sqrt[^{^{^n}}]{\frac{\sin^{2n+1}\theta}{\sqrt{n\pi}}}=\frac{1}{\sin^{2}\theta}\ge1,
\] 
因此 $\sum\limits_{n=1}^{\infty}\frac{(2n-1)!!}{(2n)!!}x^{2n} \sin ^{2n+1}\theta$ 在 $(-1,1)$ 上收敛, 从而在 $[0,t]$ 上一致收敛.

\item (1.10) 至 (1.11). 当 $t\in[0,1]$ 时, 以 $t$ 为变元的幂级数
\[
\sum_{n=1}^{\infty}\frac{(2n-1)!!}{(2n+1)(2n)!!}t^{2n+1}\sin^{2n+1}\theta\le\sum_{n=1}^{\infty}\frac{(2n-1)!!}{(2n+1)(2n)!!}=\frac{\pi}{2}-1
\]
满足 Weierstrass 判别法的条件, 因此 $\sum\limits_{n=1}^{\infty}\frac{(2n-1)!!}{(2n+1)(2n)!!}t^{2n+1}\sin^{2n+1}\theta$ 在 $[0,1]$ 上一致收敛.

\item (1.12) 至 (1.13). 同理, 根据 Weierstrass 判别法
\[
\sum_{n=1}^{\infty}\frac{(2n-1)!!}{(2n+1)(2n)!!}\sin ^{2n+1}\theta
\]
在 $\left[0,\dfrac{\pi}{2}\right]$ 上一致收敛. 
\end{enumerate}

之后的证明是相同的.







\end{document}